% concatenated from Nuclear_WCT_operators.tex
\documentclass{amsart}
\usepackage{latexsym,amsmath,amssymb}
\usepackage{cite}
\usepackage{tikz}
\usepackage{caption}
%\usepackage{tikz}
\usepackage{tikz-3dplot}
%\usepackage{caption}

%\geometry{margin=1in}
%\documentstyle[12pt]{article}
% ----------------------------------------------------------------
%\vfuzz2pt \hfuzz2pt
% THEOREMS -------------------------------------------------------
\newtheorem{thm}{Theorem}[section]
\newtheorem{cor}[thm]{Corollary}
\newtheorem{exam}[thm]{Example}
\newtheorem{lem}[thm]{Lemma}
\newtheorem{prop}[thm]{Proposition}
\theoremstyle{definition}\newtheorem{defn}[thm]{Definition}
\theoremstyle{remark}
\newtheorem{rem}[thm]{Remark}
\numberwithin{equation}{section}
% MATH -----------------------------------------------------------
\newcommand{\norm}[1]{\left\Vert#1\right\Vert}
\newcommand{\abs}[1]{\left\vert#1\right\vert}
\newcommand{\set}[1]{\left\{#1\right\}}
\newcommand{\Real}{\mathbb R}
\newcommand{\eps}{\varepsilon}
\newcommand{\To}{\longrightarrow}
\newcommand{\BX}{\mathbf{B}(X)}
\newcommand{\A}{\mathcal{A}}
\newcommand{\B}{\mathcal{B}}
\newcommand{\R}{\mathcal{R}}
\newcommand{\N}{\mathcal{N}}
\newcommand{\Complex}{\mathbb{C}}
\newcommand{\va}{\varphi}
\newcommand{\ci}{\circledcirc}
\begin{document}

\title[]
{Weighted conditional expectation operators and nuclearity}

\author{\sc\bf M. S. Al Ghafri, S. Shamsigamchi and Y. Estaremi }
\address{\sc M. S. Al Ghafri, S. Shamsigamchi and Y. Estaremi}
\email{mohammed.alghafri@utas.edu.om}\email{S.shamsi@pnu.ac.ir}\email{y.estaremi@gu.ac.ir}
\address{Department of Mathematics, University of Technology and Applied Sciences, Rustaq {329},  Oman.}
\address{Department of Mathematics, Payame Noor(PNU) university, Tehran, Iran.}
\address{Department of Mathematics and Computer Sciences,
Golestan University, Gorgan, Iran.}
%\address{\sc S. Shamsigamchi}
%\address{Department of Mathematics, Payame Noor(PNU) university, Tehran, Iran.}
%\email{S.shamsi@pnu.ac.ir}
%\address{\sc M. S. Al Ghafri}
%\email{mohammed.alghafri@utas.edu.om}
%\address{Department of Mathematics, University of Technology and Applied Sciences, Rustaq {329}, Oman.}

%\address{}
\thanks{}

\thanks{}

\subjclass[2020]{47B38}

\keywords{Weighted conditional type operator, Nuclear operator, Absolutely summing operator.}

\date{}

\dedicatory{}

\commby{}

%%% ----------------------------------------------------------------------
\begin{abstract}
We provide a characterisations of nuclear weighted conditional expectation operators on $L^p(\mu)$-spaces, for $1\leq p<\infty$. As a consequence, when the underlying measure space is non-atomic, the only nuclear weighted conditional expectation operator on $L^p(\mu)$-spaces is the zero operator.
\end{abstract}

\maketitle

\section{ \sc\bf Introduction and Preliminaries}
%Let $X$ and $Y$ be Banach spaces and $X^*$, $Y^*$ be their Duals(the set of all continuous or (equivalently) bounded linear functionals on $X$ and $Y$ with the usual norm), respectively. A linear operator $T:X\rightarrow Y$ is called nuclear if there are two sequences $\{f_n\}_{n\in \mathbb{N}}\subseteq X^*$ and $\{y_n\}_{n\in \mathbb{N}}\subseteq Y$ such that $\sum^{\infty}_{n=1}\|f_n\|_{X^*}\|y_n\|_{X}<\infty$ and
%\begin{equation}\label{A}
%Tx=\sum^{\infty}_{n=1}f_n\otimes y_n(x)=\sum^{\infty}_{n=1}f_n(x)y_n, \ \ \ \ x\in X.
%\end{equation}
%Define
%$$\|T\|_N:=\inf\{\sum^{\infty}_{n=1}\|f_n\|_{X^*}\|y_n\|_X\},$$
%in which the infimum is taken over all representations of $T$ in the form \ref{A}. The norm $\|.\|_N$ is known as nuclear norm and it's clear that if $T$ is nuclear, then it's nuclear norm is finite and $\|T\|\leq \|T\|_N$.
%
%
%It is easy to check that the operator defined as
%$$f_n\otimes y_n:X\rightarrow Y, \ \ \ f_n\otimes y_n(x)=f_n(x)y_n,$$
%is a rank one operator. Since finite rank operators are compact and the series $T=\sum^{\infty}_{n=1}f_n\otimes y_n$ is absolutely convergent, then nuclear operators are compact. It is a direct consequence from the definition of nuclear operators that if $S, T:X\rightarrow X$ are nuclear, then $S+T$ is also nuclear.\\
%If $S, T:X\rightarrow X$ are linear operators on the Banach spaces $X$ and $T$ is nuclear, the $ST$ and $TS$ are also nuclear, because
%$$TS=\sum^{\infty}_{n=1}g_n\otimes y_n, \ \ \ \ ST=\sum^{\infty}_{n=1}f_n\otimes z_n, \ \ g_n=x^*_n\circ S,  \ \ z_n=S(y_n),$$
%and also
%$$\sum^{\infty}_{n=1}\|g_n\|\|y_n\|\leq \|S\|\sum^{\infty}_{n=1}\|f_n\|\|y_n\|<\infty, \ \ \sum^{\infty}_{n=1}\|f_n\|\|z_n\|\leq \|S\|\sum^{\infty}_{n=1}\|f_n\|\|y_n\|<\infty.$$
%This implies that the set of all nuclear operators on the Banach space $X$ is a two sided ideal in the space of all linear operators on $X$.
%
%
%Also, the linear operator $T$ is said to be absolutely summing there exists $M>0$ such that
%$$\sum^{n}_{i=1}\|T(x_i)\|\leq M \sup_{f\in X^*, \|f\|\leq 1}|f(x_i)|=M\sup_{|\alpha_i|=1}\|\sum^n_{i=1}\alpha_ix_i\|,$$
%for all finite sequences $\{x_i\}^n_{i=1}\subseteq X$.\\
Let $(X,\Sigma,\mu)$ be a complete $\sigma$-finite measure space.
For any sub-$\sigma$-finite algebra $\mathcal{A}\subseteq
 \Sigma$, the $L^p$-space
$L^p(X,\mathcal{A},\mu_{\mid_{\mathcal{A}}})$ is abbreviated  by
$L^p(\mathcal{A})$ and its norm is denoted by $\|.\|_p$, $(1\leq p\leq \infty)$. All
comparisons between two functions or two sets are to be
interpreted as holding up to a $\mu$-null set. The support of a
measurable function $f$ is defined as $S(f)=\{x\in X; f(x)\neq
0\}$. We denote the vector space of all equivalence classes of
almost everywhere finite valued measurable functions on $X$ by
$L^0(\Sigma)$.

\vspace*{0.3cm} For a $\sigma$-finite subalgebra
$\mathcal{A}\subseteq\Sigma$, the conditional expectation operator
associated with $\mathcal{A}$ is the mapping $f\rightarrow
E^{\mathcal{A}}f$, defined for all non-negative measurable
functions $f$ as well as for all $f\in L^p(\Sigma)$, $(1\leq p\leq \infty)$, where
$E^{\mathcal{A}}f$, by the Radon-Nikodym theorem, is the unique
$\mathcal{A}$-measurable function satisfying
$\int_{A}fd\mu=\int_{A}E^{\mathcal{A}}fd\mu, \ \ \ \forall A\in
\mathcal{A}$. As an operator on $L^{p}({\Sigma})$,
$E^{\mathcal{A}}$ is idempotent and
$E^{\mathcal{A}}(L^p(\Sigma))=L^p(\mathcal{A})$. This operator
will play a major role in our work. Let $f\in L^0(\Sigma)$, then
$f$ is said to be conditionable with respect to $E$ if
$f\in\mathcal{D}(E):=\{g\in L^0(\Sigma): E(|g|)\in
L^0(\mathcal{A})\}$. We
write $E(f)$ in place of $E^{\mathcal{A}}(f)$.\\
Here we recall some basic properties of the conditional expectation operator $E$ on Hilbert space $L^p(\mu)$. Let $f,g\in L^p(\mu)$. Then\\
\vspace*{0.2cm} \noindent $\bullet$ \  If $g$ is
$\mathcal{A}$-measurable, then $E(fg)=E(f)g$.

\noindent $\bullet$ \ $|E(f)|^p\leq E(|f|^p)$.

\noindent $\bullet$ \ If $f\geq 0$, then $E(f)\geq 0$; if $f>0$,
then $E(f)>0$.

\noindent $\bullet$ \ $|E(fg)|\leq
E(|f|^p)|^{\frac{1}{p}}E(|g|^{p'})|^{\frac{1}{p'}}$, where
$\frac{1}{p}+\frac{1}{p'}=1$ (H\"{o}lder inequality).

\noindent $\bullet$ \ For each $f\geq 0$, $S(f)\subseteq
S(E(f))$. \\

 A detailed discussion and verification of
most of these properties may be found in \cite{rao}. We recall that
an $\mathcal{A}$-atom of the measure $\mu$ is an element
$A\in\mathcal{A}$ with $\mu(A)>0$ such that for each
$F\in\mathcal{A}$, if $F\subseteq A$, then either $\mu(F)=0$ or
$\mu(F)=\mu(A)$. A measure space $(X,\Sigma,\mu)$ with no atoms is
called non-atomic measure space. It is well-known fact that every
$\sigma$-finite measure space $(X,
\mathcal{A},\mu_{\mid_{\mathcal{A}}})$ can be partitioned uniquely
as $X=\left (\bigcup_{n\in\mathbb{N}}A_n\right )\cup B$, where
$\{A_n\}_{n\in\mathbb{N}}$ is a countable collection of pairwise
disjoint $\mathcal{A}$-atoms and $B$, being disjoint from each
$A_n$, is non-atomic (see \cite{z}).



It is easy to see that if the Sigma subalgebra $\mathcal{A}$ is generated by a sequence of measurable sets $\{A_n\}_{n\in \mathbb{N}}$ with positive measures, then for every $f\in\mathcal{D}(E)$,
$$E(f)=\sum_{n=1}^{\infty}\left(\frac{1}{\mu(A_n)}\int_{A_n}fd\mu\right).\chi_{A_n}$$

\begin{defn}
Let $(X,\Sigma,\mu)$ be a $\sigma$-finite measure space and let $\mathcal{A}$ be a
$\sigma$-subalgebra of $\Sigma$, such that $(X,\mathcal{A},\mu)$ is also $\sigma$-finite. Suppose that $E$ is the corresponding conditional
expectation operator on $L^p(\mu)$ relative to $\mathcal{A}$. If $1\leq p\leq\infty$, $w,u \in L^0(\mu)$ such that $uf$ is conditionable, for all $f\in L^p(\mu)$, then the corresponding weighted conditional expectation(WCE) operator is the linear
transformation $M_wEM_u:L^p(\mu)\rightarrow L^0(\mu)$ defined by $f\rightarrow wE(uf)$.
\end{defn}
Interested readers can find more information about WCE operators in \cite{ej,her,dhd}.\\
Many operator-theoretic properties of WCE operators on $L^p$-spaces were studied in \cite{e2,e,e1,ej,her,dhd} and such operators acting between distinct $L^p$-spaces in \cite{ej, ej2}. In this paper we characterize nuclear WCE operators on $L^p(\mu)$-spaces, for $1\leq p<\infty$.

\section{ \sc\bf Main Results}
Let $X$ and $Y$ be Banach spaces and $X^*$ be the dual of $X$.
A bounded operator $T:X\to Y$ is \emph{nuclear} if it admits a representation
\begin{equation}\label{A}
Tx=\sum_{n=1}^{\infty} f_n(x)\,y_n,
\qquad x\in X,
\end{equation}
for some $(f_n)\subset X^*$ and $(y_n)\subset Y$ with
$\sum_{n=1}^{\infty}\|f_n\|_{X^*}\,\|y_n\|_{Y}<\infty$.
The \emph{nuclear norm} is
\[
\|T\|_{N}
:=\inf\left\{\sum_{n=1}^{\infty}\|f_n\|_{X^*}\,\|y_n\|_{Y}:\ T=\sum_{n=1}^{\infty} f_n\otimes y_n\right\},
\]
where $(f\otimes y)(x):=f(x)\,y$. We refer to \cite{pi,die} for background on nuclear and summing operators,
including the ideal properties and standard consequences used implicitly below.

\medskip
An operator $T:X\to Y$ is \emph{absolutely $1$--summing} if there exists $M>0$ such that for every finite family
$x_1,\dots,x_m\in X$,
\begin{equation}\label{eq:1summing}
\sum_{i=1}^{m}\|T(x_i)\|
\le M\,\sup_{f\in X^*,\,\|f\|\le 1}\sum_{i=1}^{m}|f(x_i)|.
\end{equation}
The least such $M$ is denoted by $\pi_1(T)$. Using the weak $\ell_1$--norm identification, \eqref{eq:1summing} is equivalent to
\[
\sum_{i=1}^{m}\|T(x_i)\|
\le M\,\sup_{|\alpha_i|=1}\left\|\sum_{i=1}^{m}\alpha_i x_i\right\|.
\]

In the following we recall the Pietsch's domination theorem that will be used in the proof of our main results.

\begin{thm}[Pietsch domination theorem {\cite[Theorem~2.12]{die}}]\label{thm:Pietsch}
Let $1\le p<\infty$ and let $T:X\to Y$ be absolutely $p$--summing.
Then there exist $C>0$ and a Borel probability measure $\mu$ on the weak$^{*}$--compact unit ball $B_{X^*}$ such that
\[
\|Tx\|
\le C\left(\int_{B_{X^*}} |x^*(x)|^p\,d\mu(x^*)\right)^{1/p},
\qquad x\in X.
\]
In particular, for $p=1$,
\[
\|Tx\|\le C\int_{B_{X^*}} |x^*(x)|\,d\mu(x^*).
\]
\end{thm}

As we have proved in \cite{ej}, the WCE operator $T=M_wEM_u$ is a bounded operator on $L^{p}(\Sigma)$ if and only if
$(E|w|^{p})^{\frac{1}{p}}(E|u|^{p'})^{\frac{1}{p'}} \in
L^{\infty}(\mathcal{A})$,
 and in this case its norm is as follows
$$\|T\|=\|(E(|w|^{p}))^{\frac{1}{p}}(E(|u|^{p'}))^{\frac{1}{p'}}\|_{\infty},$$
in which $1<p<\infty$ and $\frac{1}{p}+\frac{1}{p'}=1$.\\
Also, the WCE operator $T$ is bounded on $L^{1}(\Sigma)$ if and only if
$uE(|w|) \in L^{\infty}(\Sigma)$ and in this case
$\|T\|=\|uE(|w|)\|_{\infty}$.\\

Moreover, in \cite{ej} it is shown that for $1<p<\infty$, WCE operator $M_wEM_u:L^{p}(\Sigma)\rightarrow
L^{p}(\Sigma)$ is compact if and only if for each $\varepsilon>0$
the set
 $$N_{\varepsilon}=\{x\in
X:(E(|u|^{p'}))^{\frac{1}{p'}}(x)(E(|w|^{p}))^{\frac{1}{p}}(x)\geq\varepsilon\},$$
consists of finitely many $\mathcal{A}$-atoms. And also, it is proved that the WCE operator $T=M_wEM_u:L^{1}(\Sigma)\rightarrow L^{1}(\Sigma)$ is
compact if and only if for each $\varepsilon>0$ the set
 $$K_\varepsilon:=\{x\in X:u(x)E(|w|)(x)\geq\varepsilon\}$$
 consists of finitely many
$\Sigma$-atoms.



\begin{cor}\label{c2.2} Let $1\leq p<\infty$. If the
measure space $(X,\mathcal{A},\mu)$ is non-atomic, then the WCE operator $T=M_wEM_u:L^{p}(\Sigma)\rightarrow L^{p}(\Sigma)$ is
nuclear if and only if  $T=0$.
\end{cor}
\begin{proof} Let $T=M_wEM_u$ be a nonzero
nuclear WCE operator. Then $T$ is compact and for the case $1<p<\infty$,
$$0\neq\|T\|=\|(E(|w|^{p}))^{\frac{1}{p}}(E(|u|^{p'}))^{\frac{1}{p'}}\|_{\infty},$$
and so there exists $\delta>0$ such that the set
$$\{x\in
X:(E(|w|^{p}))^{\frac{1}{p}}(x)(E(|u|^{p'}))^{\frac{1}{p'}}(x)>\delta\},$$
contains an $\mathcal{A}$-measurable subset $B$ with positive
measure. Hence, we obtain a sequence of pairwise disjoint subsets
$B_{n}\subseteq B$ such that for every $n\in\mathbb{N}$,
$B_{n}\in\A$ and $0<\mu(B_{n})<\infty$. Define
$$f_{n}=\frac{\overline{u}|u|^{p'-2}(E(|w|^{p}))^{\frac{p'-1}{p}}}
{\|T\|^{\frac{p'}{p}} (\mu(B_{n}))^{\frac{1}{p}}}\chi_{B_{n}}.$$
It is easy to see that $\|f_{n}\|_{p}\leq1$ and
\begin{align*}
\|Tf_{n}-Tf_{m}\|^{p}_{p}&=\int_{X}|w|^{p}|E(f_{n}-f_{m})|^{p}d\mu\\
&\geq\int_{B_{n}}\frac{(E(|w|^{p}))^{p'}(E(|u|^{p'}))^{p}}{\|T\|^{p'}\mu(B_{n})}d\mu\\
&+ \int_{B_{m}}\frac{(E(|w|^{p}))^{p'}(E(|u|^{p'}))^{p}}{\|T\|^{p'}\mu(B_{m})}d\mu\\
 &\geq \frac{2\delta^{pp'}} {\|T\|^{p'}}.
 \end{align*}
 Also, for the case $p=1$, we have
 $0\neq\|T\|=\|uE(|w|)\|_{\infty}$ and so $\mu(S(uE(|w|))>0$. Also, we have $S(uE(|w|)\subseteq S(E(|u|)E(|w|)$ and $\|E(|u|)E(|w|)\|_{\infty}\leq \|uE(|w|)\|_{\infty}<\infty$. Hence $\mu(S(E(|u|)E(|w|))>0$ and then there exists $\delta>0$ such that the set
$$\{x\in
X:E(|u|)(x)E(|w|)(x)>\delta\},$$
contains an $\mathcal{A}$-measurable subset $B$ with positive
measure. Hence, we obtain a sequence of pairwise disjoint subsets
$B_{n}\subseteq B$ such that for every $n\in\mathbb{N}$,
$B_{n}\in\A$ and $0<\mu(B_{n})<\infty$. Define
$$f_{n}=\frac{\overline{E(u)}E(|w|)}
{(\|T\|^2+1)\mu(B_{n})}\chi_{B_{n}}.$$
It is easy to see that $\|f_{n}\|_{1}\leq1$ and $\{Tf_n\}_{n\in \mathbb{N}}$ has no convergent subsequence.

  This implies that in both cases, $T$ is not compact. But this is a contradiction. The converse is obvious.\\
\end{proof}

\begin{cor}\label{c2.4}
If the measure space $(X, \mathcal{A}, \mu)$, with finite subset property, is not purely atomic, such that the restriction of the WCE operator $T=M_wEM_u$ to the non-atomic part is a non-zero operator, then $T$ can not be nuclear on $L^p(\Sigma)$ for $1\leq p<\infty$.
\end{cor}
\begin{proof}
Similar to the proof of Corollary \ref{c2.2}, if there is a non-atomic measurable set $A\in \mathcal{A}$, with $\mu(A)>0$, then we can find a bounded sequence $\{f_n\}$ in
$L^p(\mathcal{A})$ such that the sequence $\{Tf_n\}$ has no convergent subsequence. So $T$ is not compact and therefore $T$ is not nuclear on $L^p(\Sigma)$.
\end{proof}

\begin{prop}\label{p2.5}
Let $1<p<\infty$ and $\mathcal{A}\subseteq \Sigma$ be a $\sigma$- finite subalgebra such that the measure space $(X, \mathcal{A}, \mu\mid_{\mathcal{A}})$ is a purely atomic measure space. Then the bounded WCE operator $T=M_wEM_u$ is nuclear on $L^p(\Sigma)$ if and only if
$$\sum^{\infty}_{n=1}(E(|w|^{p})(A_n))^{\frac{1}{p}}(E(|u|^{q})(A_n))^{\frac{1}{q}}<\infty,$$
in which $\{A_n\}^{\infty}_{n=1}$ pair-wise disjoint $\mathcal{A}$-atoms such that $X=\cup^{\infty}_{n=0}A_n$ and $\frac{1}{p}+\frac{1}{q}=1$.
\end{prop}
\begin{proof} Since the measure space $(X, \mathcal{A}, \mu\mid_{\mathcal{A}})$ is a $\sigma$-finite and purely atomic, then there exists a sequence of pair-wise disjoint $\mathcal{A}$-atoms $\{A_n\}^{\infty}_{n=1}$ such that $X=\cup^{\infty}_{n=0}A_n$. As is known in the literature the $\mathcal{A}$-measurable functions defined on $X$ are fixed on each $\mathcal{A}$-atom. Hence for every $f\in L^p(\Sigma)$ we have
\begin{align*}
Tf=w.E(uf)&=w.\sum_{n=1}^{\infty}E(u.f)(A_n).\chi_{A_n}\\
&=w.\sum_{n=1}^{\infty}\left(\frac{1}{\mu(A_n)}\int_{A_n}E(uf)d\mu\right).\chi_{A_n}\\
&=w.\sum_{n=1}^{\infty}\left(\frac{1}{\mu(A_n)}\int_{A_n}ufd\mu\right).\chi_{A_n}.
\end{align*}
Let $$\sum^{\infty}_{n=1}(E(|w|^{p})(A_n))^{\frac{1}{p}}(E(|u|^{q})(A_n))^{\frac{1}{q}}<\infty.$$
 Since $T$ is bounded and the measure space $(X, \mathcal{A}, \mu\mid_{\mathcal{A}})$ is a purely atomic, then we have
$$\|T\|=\|(E(|w|^{p}))^{\frac{1}{p}}(E(|u|^{q}))^{\frac{1}{q}}\|_{\infty}=\sup_{n}(E(|w|^{p})(A_n))^{\frac{1}{p}}(E(|u|^{q})(A_n))^{\frac{1}{q}}<\infty.$$
If we set $\phi_n(f)=\int_X(u.\chi_{A_n})fd\mu$ and $g_n=\frac{w.\chi_{A_n}}{\mu(A_n)}$, then we get that for every $n\in \mathbb{N}$, $u.\chi_{A_n}\in L^{q}(\Sigma)$ and  $g_n\in L^{p}(\Sigma)$. It is clear that for each $n\in \mathbb{N}$, the function $\phi_n$ is a bounded linear functional on $L^p(\Sigma)$. Also, $\|g_n\|_p=(E(|w|^p)(A_n))^{\frac{1}{p}}$ and $\|\phi_n\|=(E(|u|^{q})(A_n))^{\frac{1}{q}}$. Therefore,

$$T=M_wEM_u=\sum_{n=1}^{\infty}\phi_n\otimes g_n,$$
and
$$\sum_{n=1}^{\infty}\|\phi_n\|\|g_n\|=\sum^{\infty}_{n=1}(E(|w|^{p})(A_n))^{\frac{1}{p}}(E(|u|^{q})(A_n))^{\frac{1}{q}}<\infty.$$
This implies that $T$ is nuclear.\\

Conversely, let the WCE operator $T=M_wEM_u$ be nuclear. Then $T$ is absolutely summing and so by Pietsch’s  theorem 2.12 of \cite{die}, there is a Borel probability measure $\nu$ on the unit ball $B_q$ of $L^q(\mu)\simeq(L^p(\mu))^*$ such that for some $M>0$

$$\|Tf\|_p\leq M\int_{B_q}|l(f)|d\nu(l), \ \ \ \ \ \forall f\in L^p(\mu).$$
By Riesz representation theorem, for each $l\in (L^p(\mu))^*$, there exists $g_l\in L^q(\mu)$ such that
  $$l(f)=\int_Xf.g_l d\mu, \ \ \ \ \ f\in L^p(\mu).$$
  
So by Holder's inequality we have 
\begin{align*}
\|Tf\|_p&\leq M\int_{B_q}|l(f)|d\nu(l)\\
&=M\int_{B_q}\int_X|f.g_l| d\mu d\nu(g_l)\\
&\leq M\int_{B_q}\|g_l\|_q\|f\|_pd\nu(g_l),
\end{align*}
 
Since $\{A_n\}^{\infty}_{n=1}$ is a sequence of disjoint atoms, for each $i\in \mathbb{N}$ we define
$$f_i=\frac{\overline{u}|u|^{q-2}(E(|w|^{p}))^{\frac{q-1}{p}}}{\|T\|^{\frac{q}{p}} (\mu(A_{i}))^{\frac{1}{p}}}\chi_{A_{i}}.$$
It is easy to see that $f_i\in L^p(\Sigma)$ and $\|f_i\|_p\leq 1$. Hence we have
\begin{align*}
\|Tf_i\|^p_p&=\frac{1}{\|T\|^{q}}\int_{A_i}\frac{1}{\mu(A_{i})}\left((E(|w|^{p}))^{\frac{1}{p}}(E(|u|^{q}))^{\frac{1}{q}}\right)^{pq}d\mu\\
&=\frac{1}{\|T\|^{q}}\left((E(|w|^{p})(A_i))^{\frac{1}{p}}(E(|u|^{q})(A_i))^{\frac{1}{q}}\right)^{pq}.
\end{align*}
By these observations we get 
\begin{align*}
\|Tf_i\|_p&=\frac{1}{\|T\|^{\frac{q}{p}}}\left((E(|w|^{p})(A_i))^{\frac{1}{p}}(E(|u|^{q})(A_i))^{\frac{1}{q}}\right)^{q}\\
&\leq M\int_{B_q}\|g_l.\chi_{A_i}\|_q\|f_i\|_pd\nu(g_l)\\
&\leq  M\int_{B_q}\|g_l.\chi_{A_i}\|_qd\nu(g_l).\\
\end{align*}
And so 
$$(E(|w|^{p})(A_i))^{\frac{1}{p}}(E(|u|^{q})(A_i))^{\frac{1}{q}}\leq M^{\frac{1}{q}}\|T\|^{\frac{1}{p}}\left(\int_{B_q}\|g_l.\chi_{A_i}\|_qd\nu(g_l)\right)^{\frac{1}{q}}.$$

It is clear that $\|g_l\|_q=\sum^{\infty}_{i=1}\|g_l.\chi_{A_i}\|_q$. Therefore
\begin{align*}
\sum^{\infty}_{i=1}(E(|w|^{p})(A_i))^{\frac{1}{p}}(E(|u|^{q})(A_i))^{\frac{1}{q}}&\leq M^{\frac{1}{q}}\|T\|^{\frac{1}{p}}\sum^{\infty}_{i=1}\left(\int_{B_q}\|g_l.\chi_{A_i}\|_qd\nu(g_l)\right)^{\frac{1}{q}}\\
&\leq M^{\frac{1}{q}}\|T\|^{\frac{1}{p}}\left(\int_{B_q}\sum^{\infty}_{i=1}\|g_l.\chi_{A_i}\|_qd\nu(g_l)\right)^{\frac{1}{q}}\\
&=M^{\frac{1}{q}}\|T\|^{\frac{1}{p}}\left(\int_{B_q}\|g_l\|_qd\nu(g_l)\right)^{\frac{1}{q}}\\
&\leq M^{\frac{1}{q}}\|T\|^{\frac{1}{p}}.
\end{align*}
This implies that 

$$\sum^{\infty}_{i=1}(E(|w|^{p})(A_i))^{\frac{1}{p}}(E(|u|^{q})(A_i))^{\frac{1}{q}}\leq M^{\frac{1}{q}}\|T\|^{\frac{1}{p}}<\infty.$$

This completes the proof.
\end{proof}

In the following theorem we characterize compact WCE operators on $L^1(\Sigma)$.
\begin{thm}
The WCE operator $T:L^{1}(\Sigma)\rightarrow L^{1}(\Sigma)$ is
compact if and only if for each $\varepsilon>0$ the set
$$K_\varepsilon:=\{x\in X:E(u)(x)E(|w|)(x)\geq\varepsilon\}$$
consists of finitely many $\mathcal{A}$-atoms.
\end{thm}
\begin{proof}
Suppose that $T$ is
compact but for some
$\varepsilon>0$ the set $K_\varepsilon$ either contains a
 subset of non-atomic part $A$ with positive measure or has infinitely many $\mathcal{A}$-atoms. In both cases we
 can find a sequence of pairwise disjoint measurable subsets
 $\{A_{n}\}_{n\in \mathbb{N}}$ with $0<\mu(A_{n})<\infty$.
 Define $f_{n}=\frac{\overline{E(u)}E(|w|)\chi_{A_{n}}}{\mu(A_{n})\|T\|^2}$. Then
 $\|f_{n}\|_{1}\leq1/\|T\|$ and $\sigma(f_n)\cap\sigma(f_m)=\emptyset$, for $n\neq
 m$. Hence we have
  Then
  \begin{align*}
  \|Tf_{n}-Tf_{m}\|_{1}&\geq\int_{A_{n}}E(|w|)|E(uf_{n}-uf_{m})|d\mu\\
 &=\int_{A_{n}}\frac{(E(|w|))^2|E(u)|^2\chi_{A_{n}}}{\mu(A_{n})\|T\|^2}d\mu\\
 &\geq\frac{\varepsilon^{2}}{\|T\|^2}.
 \end{align*}
which shows that the sequence $\{Tf_n\}$ dose not contain a
convergent subsequence, and so $T$ is not compact.\\

Conversely, suppose that for each $\varepsilon>0$,
$K_\varepsilon=\cup^{n}_ {k=1}A^{\varepsilon}_{k}$. Define
$T_{\varepsilon}(f)=T(f\chi_{K_\varepsilon})=\Sigma_{k=1}^nT(f\chi_{A^{\varepsilon}_{k}})$,
for all $f\in L^{1}(\Sigma)$. It is clear that $T_{\varepsilon}$ is
a finitely rank operator on $L^{1}(\Sigma)$ and
$\|T-T_{\varepsilon}\|<\varepsilon$. Thus $T$ is compact.
\end{proof}

\begin{prop}\label{p2.7}
Let $\mathcal{A}\subseteq \Sigma$ be a $\sigma$- finite subalgebra such that the measure space $(X, \mathcal{A}, \mu)$ is a purely atomic measure space. Then the bounded WCE operator $T=M_wEM_u:L^1(\Sigma)\rightarrow L^1(\Sigma)$ is nuclear on if and only if
$$\sum^{\infty}_{n=1}E(|w|)(A_n)E(|u|)(A_n)<\infty,$$
in which $\{A_n\}^{\infty}_{n=1}$ pair-wise disjoint $\mathcal{A}$-atoms such that $X=\cup^{\infty}_{n=0}A_n$.
\end{prop}
\begin{proof}
Slimilar to the proof of Proposition \ref{p2.5} for every $f\in L^p(\Sigma)$ we have
\begin{align*}
Tf=w.E(uf)&=w.\sum_{n=1}^{\infty}E(u.f)(A_n).\chi_{A_n}\\
&=w.\sum_{n=1}^{\infty}\left(\frac{1}{\mu(A_n)}\int_{A_n}E(uf)d\mu\right).\chi_{A_n}\\
&=w.\sum_{n=1}^{\infty}\left(\frac{1}{\mu(A_n)}\int_{A_n}ufd\mu\right).\chi_{A_n}.
\end{align*}
Let
$$\sum^{\infty}_{n=1}E(|w|)(A_n)E(|u|)(A_n)<\infty.$$
 Since $T$ is bounded and the measure space $(X, \mathcal{A}, \mu\mid_{\mathcal{A}})$ is a purely atomic, then we have
$$\sup_{n\in \mathbb{N}}\left[E(|w|)(A_n)E(|u|)(A_n)\right]=\|E(|w|)E(|u|)\|_{\infty}\leq\|uE(|w|)\|_{\infty}=\|T\|.$$
If we set $\phi_n(f)=\int_X(u.\chi_{A_n})fd\mu$ and $g_n=\frac{w.\chi_{A_n}}{\mu(A_n)}$, then we get that for every $n\in \mathbb{N}$, $u.\chi_{A_n}\in L^{\infty}(\Sigma)$ and  $g_n\in L^{1}(\Sigma)$. It is clear that for each $n\in \mathbb{N}$, the function $\phi_n$ is a bounded linear functional on $L^1(\Sigma)$. Also, $\|g_n\|_1=E(|w|)(A_n)$ and $\|\phi_n\|\leq E(|u|)(A_n)$. Therefore,

$$T=M_wEM_u=\sum_{n=1}^{\infty}\phi_n\otimes g_n,$$
and
$$\sum_{n=1}^{\infty}\|\phi_n\|\|g_n\|\leq\sum^{\infty}_{n=1}E(|w|)(A_n)E(|u|)(A_n)<\infty.$$
This implies that $T$ is nuclear.\\

Conversely, let the WCE operator $T=M_wEM_u$ be nuclear on $L^1(\mu)$. Then $T$ is absolutely summing and so by Pietsch’s  theorem 2.12 of \cite{die}, there is a Borel probability measure $\nu$ on the unit ball $B_{\infty}$ of $L^{\infty}(\mu)\simeq(L^1(\mu))^*$ such that for some $M>0$

$$\|Tf\|_1\leq M\int_{B_{\infty}}|l(f)|d\nu(l), \ \ \ \ \ \forall f\in L^1(\mu).$$
By Riesz representation theorem, for each $l\in (L^1(\mu))^*$, there exists $g_l\in L^{\infty}(\mu)$ such that
  $$l(f)=\int_Xf.g_l d\mu, \ \ \ \ \ f\in L^1(\mu).$$
  
And hence we have
\begin{align*}
\|Tf\|_1&\leq M\int_{B_{\infty}}|l(f)|d\nu(l)\\
&=M\int_{B_{\infty}}\int_X|f.g_l| d\mu d\nu(g_l)\\
&\leq M\int_{B_{\infty}}\|g_l\|_{\infty}\|f\|_1d\nu(g_l),
\end{align*}
 
Since $\{A_n\}^{\infty}_{n=1}$ is a sequence of disjoint atoms, for each $n\in \mathbb{N}$ we define
$$f_{n}=\frac{\overline{E(u)}E(|w|)\chi_{A_{n}}}{\mu(A_{n})\|T\|(\|T\|+1)}.$$
It is easy to see that $f_n\in L^1(\Sigma)$ and $\|f_n\|_1\leq 1$. Hence we have

\begin{align*}
\|Tf_n\|_1&=\frac{1}{\|T\|(\|T\|+1)}\int_{A_n}\frac{1}{\mu(A_{n})}\left(E(|w|)E(|u|)\right)^{2}d\mu\\
&=\frac{1}{\|T\|(\|T\|+1)}\left(E(|w|)(A_n)E(|u|)(A_n)\right)^{2}.
\end{align*}
By these observations we get that
\begin{align*}
\|Tf_n\|_1&=\frac{1}{\|T\|(\|T\|+1)}\left(E(|w|)(A_n)E(|u|)(A_n)\right)^{2}\\
&\leq M\int_{B_{\infty}}\|g_l.\chi_{A_n}\|_{\infty}\|f\|_1d\nu(g_l),\\
&\leq  M\int_{B_{\infty}}\|g_l.\chi_{A_n}\|_{\infty}d\nu(g_l).\\
\end{align*}
And so
$$E(|w|)(A_n)E(|u|)(A_n)\leq  M^{\frac{1}{2}}(\|T\|(\|T\|+1))^{\frac{1}{2}}\left(\int_{B_{\infty}}\|g_l.\chi_{A_n}\|_{\infty}d\nu(g_l)\right)^{\frac{1}{2}}.\\$$
It is clear that $\|g_l\|_{\infty}=\sum^{\infty}_{n=1}\|g_l.\chi_{A_n}\|_{\infty}$. Therefore
\begin{align*}
\sum^{\infty}_{i=1}E(|w|)(A_n)E(|u|)(A_n)&\leq M^{\frac{1}{2}}(\|T\|(\|T\|+1))^{\frac{1}{2}}\sum^{\infty}_{i=1}\left(\int_{B_{\infty}}\|g_l.\chi_{A_n}\|_{\infty}d\nu(g_l)\right)^{\frac{1}{2}}\\
&\leq M^{\frac{1}{2}}(\|T\|(\|T\|+1))^{\frac{1}{2}}\left(\int_{B_{\infty}}\sum^{\infty}_{i=1}\|g_l.\chi_{A_n}\|_{\infty}d\nu(g_l)\right)^{\frac{1}{2}}\\
&=M^{\frac{1}{2}}(\|T\|(\|T\|+1))^{\frac{1}{2}}\left(\int_{B_{\infty}}\|g_l\|_{\infty}d\nu(g_l)\right)^{\frac{1}{2}}\\
&\leq M^{\frac{1}{2}}(\|T\|(\|T\|+1))^{\frac{1}{2}}.
\end{align*}
This implies that 
$$\sum^{\infty}_{i=1}E(|w|)(A_n)E(|u|)(A_n)\leq M^{\frac{1}{2}}(\|T\|(\|T\|+1))^{\frac{1}{2}}<\infty.$$
\end{proof}

\begin{rem}
The WCE operator $T=M_wEM_u:L^p(\mu)\rightarrow L^p(\mu)$ is nuclear if and only if $T^*$ is nuclear, for all $1\leq p<\infty$.
\end{rem}

Let $\mathcal{A}\subseteq \Sigma$ be a $\sigma$-finite subalgebra of $\Sigma$. Then $X=\left (\bigcup_{n\in\mathbb{N}}A_n\right )\cup B$, where
$\{A_n\}_{n\in\mathbb{N}}$ is a countable collection of pairwise disjoint $\mathcal{A}$-atoms and $B$, being disjoint from each
$A_n$, is non-atomic. In the following theorem we characterize nuclear WCE operators on general $L^p(\mu)$-spaces.

  \begin{thm}\label{t2.10}
Let $1<p<\infty$ and $\mathcal{A}\subseteq \Sigma$ be a $\sigma$- finite subalgebra. Then the bounded WCE operator $T=M_wEM_u$ is nuclear on $L^p(\Sigma)$ if and only if
\begin{itemize}
  \item $(E(|w|^{p}))^{\frac{1}{p}}(E(|u|^{q})))^{\frac{1}{q}}=0$, $\mu$-a.e., on $B$. 
  \item $\sum^{\infty}_{n=1}(E(|w|^{p})(A_n))^{\frac{1}{p}}(E(|u|^{q})(A_n))^{\frac{1}{q}}<\infty$,
\end{itemize}
in which $\{A_n\}^{\infty}_{n=1}$ pair-wise disjoint $\mathcal{A}$-atoms and $B$ is non-atomic such that $X=\cup^{\infty}_{n=0}A_n\cup B$ and $\frac{1}{p}+\frac{1}{q}=1$.
\end{thm}
\begin{proof}
By combining Proposition \ref{p2.5} and Corollary \ref{c2.4} we have the result.
\end{proof}
\begin{thm}\label{t2.11}
Let $\mathcal{A}\subseteq \Sigma$ be a $\sigma$- finite subalgebra. Then the bounded WCE operator $T=M_wEM_u$ is nuclear on $L^1(\Sigma)$ if and only if
\begin{itemize}
  \item $E(|w|)E(|u|)=0$, $\mu$-a.e., on $B$.
  \item $\sum^{\infty}_{n=1}E(|w|)(A_n)E(|u|)(A_n)<\infty$,
\end{itemize}
in which $\{A_n\}^{\infty}_{n=1}$ pair-wise disjoint $\mathcal{A}$-atoms and $B$ is non-atomic such that $X=\cup^{\infty}_{n=0}A_n\cup B$ and $\frac{1}{p}+\frac{1}{q}=1$.
\end{thm}
\begin{proof}
By combining Proposition \ref{p2.7} and Corollary \ref{c2.4} we have the result.
\end{proof}
\begin{cor}
Let $\mathcal{A}\subseteq \Sigma$ be a $\sigma$- finite subalgebra and $1\leq p<\infty$. Then the multiplication operator $M_u:L^p(\mu)\rightarrow L^p(\mu)$ is nuclear if and only if 
$u=0$, $\mu$-a.e., on $B$ and 
$$\sum^{\infty}_{n=1}|u|(A_n)<\infty.$$

\end{cor}
\begin{exam}
Let $X=\mathbb{N}$,
$\Sigma=2^{\mathbb{N}}$ and let $\mu$ be the counting measure. Put
$$A=\{\{2\},\{4,6\},
\{8,10,12\},\{14,16,18,20\},\cdots\}\cup\{\{1\}, \{3\}, \{5\},
\cdots \}.$$
If we let $A_{1}=\{2\}$, $A_{2}=\{4,6\}$,
$A_{3}=\{8,10,12\}$, $\cdots$, then we see that $\mu(A_{n})=n$
and for every $n\in \mathbb{N}$, there exists $k_{n}\in
\mathbb{N}$ such that $A_{n}=\{2k_{n}, 2(k_{n}+1), \cdots,
2(k_{n}+n-1)\}$. Let $\mathcal{A}$ be the $\sigma$-algebra
generated by the partition $A$ of $\mathbb{N}$. Note that,
$\mathcal{A}$ is a sub-$\sigma$-finite algebra of $\Sigma$ and
each of element of $\mathcal{A}$ is an $\mathcal{A}$-atom. It is
known that the conditional expectation of any  $f\in
\mathcal{D}(E)$ relative to $\mathcal{A}$ is
$$E(f)=\sum_{n=1}^{\infty}\left(\frac{1}{\mu(A_n)}\int_{A_n}fd\mu\right)\chi_{A_n}+
\sum_{n=1}^{\infty}f(2n-1)\chi_{\{2n-1\}}.$$
Define $u(n)=n$ and $w(n)=\frac{1}{n^{3}}$, for all $n\in \mathbb{N}$.
For each even number $m\in \mathbb{N}$, there exists $n_{m}\in \mathbb{N}$ such that
$m\in A_{n_{m}}$. Hence for every$ 1< p<\infty$ we have
$$E(|w|^p)(m)=\frac{1}{(2k_{n_m})^{3p}}+\cdots
+\frac{1}{(2k_{n_m}+2n_m-2)^{3p}}$$
and
$$E(|u|^{q})(m)=2^qk^q_{n_m}+2^q(k_{n_m}+1)^q+ \cdots+
2^q(k_{n_m}+n_m-1)^q$$
Since $n_m\leq k_{n_m}$ we have

$$(E(|w|^p))^{\frac{1}{p}}(m)(E(|u|^{q}))^{\frac{1}{q}}(m)\leq\frac{4k_{n_m}}{2^3k^3_{n_m}}.$$

Also, for all $n\in \mathbb{N}$
$$(E(|w|^p))^{\frac{1}{p}}(2n-1)(E(|u|^{p}))^{\frac{1}{p}}(2n-1)\leq
\frac{1}{(2n-1)^2}.$$
Therefore 
\begin{align*}
\sum^{\infty}_{n=1}(E(|w|^{p})(A_n))^{\frac{1}{p}}(E(|u|^{q})(A_n))^{\frac{1}{q}}&+\sum^{\infty}_{n=1}(E(|w|^{p})(2n-1))^{\frac{1}{p}}(E(|u|^{q})(2n-1))^{\frac{1}{q}}\\
&\leq \sum^{\infty}_{n=1}\frac{1}{2k^2_{n}}+\sum^{\infty}_{n=1}\frac{1}{(2n-1)^2}\\
&<\infty,
\end{align*}
and so by Theorem \ref{t2.10} we get that the WCT operator $T=M_wEM_u$ is a nuclear operator on $L^p(\mu)$. Similarly we have
\begin{align*}
\sum^{\infty}_{n=1}E(|w|)(A_n)E(|u|)(A_n)&+\sum^{\infty}_{n=1}E(|w|)(2n-1)E(|u|)(2n-1)\\
&\leq \sum^{\infty}_{n=1}\frac{1}{2k^2_{n}}+\sum^{\infty}_{n=1}\frac{1}{(2n-1)^2}\\
&<\infty,
\end{align*}
and consequently by Theorem \ref{t2.11} we get that we get that the WCT operator $T=M_wEM_u$ is a nuclear operator on $L^1(\mu)$.

\end{exam}















%
%
%\section*{1. Definition of $\sigma(g)$}
%
%Let $g : (\Omega,\mathcal{F}) \to (S,\mathcal{B})$ be a measurable function.
%The sigma-algebra generated by $g$ is defined as
%\[
%\sigma(g) = \{\, g^{-1}(B) : B \in \mathcal{B} \,\}.
%\]
%
%This is the smallest $\sigma$-algebra on $\Omega$ making $g$ measurable.
%
%\bigskip
%
%\noindent
%\textbf{Interpretation:}
%A set $A \subseteq \Omega$ belongs to $\sigma(g)$ iff its membership can be
%determined solely from the value of $g(\omega)$.
%
%\bigskip
%
%\section*{2. Characterizations of $\sigma(g)$}
%
%We present several equivalent and useful descriptions of $\sigma(g)$.
%
%\subsection*{(a) Preimage Characterization}
%
%\[
%\sigma(g) = \{ g^{-1}(B) : B\in\mathcal{B}(S) \}.
%\]
%
%\subsection*{(b) Characterization via Fibers}
%
%For each $s\in S$ define the \emph{fiber}
%\[
%g^{-1}(\{s\}) = \{\omega \in \Omega : g(\omega)=s\}.
%\]
%
%These fibers form a partition of $\Omega$.
%Then:
%\[
%A \in \sigma(g)
%\quad\Longleftrightarrow\quad
%A \text{ is a union of fibers of } g.
%\]
%
%Equivalently, for every $s$,
%\[
%A \cap g^{-1}(\{s\}) \text{ is either the entire fiber or empty}.
%\]
%
%\subsection*{(c) Measurable Functions w.r.t.\ $\sigma(g)$}
%
%A function $f: \Omega \to \mathbb{R}$ is $\sigma(g)$-measurable iff
%\[
%f = h \circ g
%\]
%for some measurable $h : S \to \mathbb{R}$.
%
%This is essential in computing conditional expectations:
%\[
%\mathbb{E}[X \mid \sigma(g)] = h(g)
%\]
%for a unique measurable $h$.
%
%\subsection*{(d) Information Content}
%
%Two points $\omega_1, \omega_2 \in \Omega$ cannot be distinguished by any
%set in $\sigma(g)$ unless $g(\omega_1) \neq g(\omega_2)$.
%
%Thus:
%\[
%\omega_1 \text{ and } \omega_2 \text{ lie in the same atom of } \sigma(g)
%\quad \Longleftrightarrow \quad
%g(\omega_1) = g(\omega_2).
%\]
%
%\section*{3. Examples}
%
%\subsection*{Example 1: $g(x) = x^2$ on $[0,1]$}
%
%Here the fibers are
%\[
%g^{-1}(y) = \{\sqrt{y},\; 1 - \sqrt{y}\}.
%\]
%
%Thus
%\[
%A \in \sigma(x^2)
%\quad\Longleftrightarrow\quad
%x \in A \iff 1-x \in A.
%\]
%
%A function $f$ is $\sigma(x^2)$-measurable iff
%\[
%f(x) = f(1-x).
%\]
%
%\subsubsection*{Conditional Expectation Example}
%
%Let $f(x)=\sin(5\pi x)$. Then
%\[
%\mathbb{E}[f \mid \sigma(x^2)](x)
%= \frac12 \left(\sin(5\pi x) + \sin(5\pi(1-x))\right).
%\]
%
%\subsection*{Example 2: $g(x)=|x|$ on $[-1,1]$}
%
%Fibers are
%\[
%g^{-1}(y) = \{y, -y\}.
%\]
%
%Thus a set $A$ belongs to $\sigma(|x|)$ iff it is symmetric:
%\[
%x \in A \iff -x \in A.
%\]
%
%A function is $\sigma(|x|)$-measurable iff $f(x)=f(-x)$.
%
%\subsection*{Example 3: $g(x)=\lfloor x \rfloor$ on $\mathbb{R}$}
%
%Fibers:
%\[
%g^{-1}(n) = [n, n+1).
%\]
%
%Thus
%\[
%\sigma(g)
%= \sigma\bigl( [n,n+1) : n \in \mathbb{Z} \bigr),
%\]
%which is the sigma-algebra consisting of unions of integer-length blocks.
%
%A function $f$ is $\sigma(g)$-measurable iff it is constant on each interval
%$[n, n+1)$.
%
%\subsection*{Example 4: Dyadic Sigma-Algebra Projection}
%
%Let
%\[
%\mathcal{G}_n = \sigma\bigl( [k2^{-n}, (k+1)2^{-n}) : k=0,\dots,2^n-1 \bigr).
%\]
%
%For $f(x)=\sin(40\pi x)$,
%\[
%\mathbb{E}[f \mid \mathcal{G}_n](x)
%= 2^n \int_{k2^{-n}}^{(k+1)2^{-n}} \sin(40\pi t)\,dt
%\]
%for $x$ in the interval $[k2^{-n},(k+1)2^{-n})$.
%
%\subsection*{Example 5: Conditional Expectation in $L^p$}
%
%Let $\Omega=[0,2]$, $\mathcal{G}=\sigma([0,1], (1,2])$ and $f(x)=x^2$. Then:
%\[
%\mathbb{E}[f \mid \mathcal{G}]
%= \frac13 \mathbf{1}_{[0,1]} + \frac{7}{3}\mathbf{1}_{(1,2]}.
%\]
%
%This works in any $L^p$ space.
%
%\section*{4. Summary}
%
%For any measurable function $g$:
%\begin{itemize}
%    \item $\sigma(g)$ consists exactly of preimages of measurable sets.
%    \item $\sigma(g)$ partitions $\Omega$ into the fibers of $g$.
%    \item A function is $\sigma(g)$-measurable iff it is a measurable transform of $g$.
%    \item Conditional expectation onto $\sigma(g)$ produces functions of the form $h(g)$.
%\end{itemize}
%
%
%
%
%
%
%\section{Rigorous Characterization of $\sigma(g)$}
%
%Let $(\Omega,\mathcal{F},\mu)$ be a measure space, and let
%$g:\Omega\to S$ be measurable into $(S,\mathcal{B})$.
%
%\begin{defn}[Sigma-Algebra Generated by $g$]
%\[
%\sigma(g) = \{\, g^{-1}(B) : B \in \mathcal{B} \,\}.
%\]
%\end{defn}
%
%\subsection{Equivalence with Fiber Partitions}
%
%\begin{thm}
%A set $A \subseteq \Omega$ belongs to $\sigma(g)$ if and only if $A$ is a union of fibers of $g$, i.e.,
%\[
%\forall s \in S, \quad g^{-1}(\{s\}) \subseteq A \text{ or } g^{-1}(\{s\}) \cap A = \emptyset.
%\]
%\end{thm}
%
%\begin{proof}
%\emph{($\Rightarrow$)} Let $A = g^{-1}(B)$ for some $B\in \mathcal{B}$. Then for each $s\in S$:
%\[
%g^{-1}(\{s\}) \subseteq A \iff s \in B, \quad g^{-1}(\{s\}) \cap A = \emptyset \iff s \notin B.
%\]
%Thus $A$ is a union of fibers.
%
%\emph{($\Leftarrow$)} Let $A$ be a union of fibers: $A = \bigcup_{s\in C} g^{-1}(\{s\})$ for some $C \subseteq S$. Since $g$ is measurable, $g^{-1}(C) = A \in \sigma(g)$.
%\end{proof}
%
%\subsection{Equivalence with Measurable Functions of $g$}
%
%\begin{thm}
%A function $f:\Omega\to \mathbb{R}$ is $\sigma(g)$-measurable iff there exists a measurable $h:S\to \mathbb{R}$ such that
%\[
%f = h \circ g.
%\]
%\end{thm}
%
%\begin{proof}
%\emph{($\Rightarrow$)} Let $f$ be $\sigma(g)$-measurable. Then for all Borel $B \subseteq \mathbb{R}$, $f^{-1}(B) \in \sigma(g)$. By the fiber characterization, $f$ is constant on each fiber $g^{-1}(\{s\})$, so define $h(s) := f(\omega)$ for any $\omega \in g^{-1}(\{s\})$.
%
%\emph{($\Leftarrow$)} If $f = h \circ g$, then for any Borel $B$, $f^{-1}(B) = g^{-1}(h^{-1}(B)) \in \sigma(g)$.
%\end{proof}
%
%\section{Conditional Expectation via Radon--Nikodym}
%
%Let $(\Omega, \mathcal{F}, \mathbb{P})$ be a probability space and $X \in L^1(\Omega)$. Define a measure $\nu$ on $\sigma(g)$:
%\[
%\nu(A) = \int_A X \, d\mathbb{P}, \quad A \in \sigma(g).
%\]
%
%By the Radon--Nikodym theorem, there exists a unique $\sigma(g)$-measurable $Y$ such that
%\[
%\nu(A) = \int_A Y \, d\mathbb{P}, \quad \forall A \in \sigma(g),
%\]
%and we define $Y = \mathbb{E}[X \mid \sigma(g)]$.
%
%Since $Y$ is $\sigma(g)$-measurable, there exists $h$ such that
%\[
%\mathbb{E}[X \mid \sigma(g)] = h \circ g.
%\]
%
%\section{Geometric Interpretation in $L^2$}
%
%Let $X \in L^2(\Omega,\mathcal{F},\mathbb{P})$. The closed subspace of $\sigma(g)$-measurable functions is
%\[
%L^2(\sigma(g)) = \{ h \circ g : h \in L^2(S, \mathcal{B}, g_\# \mathbb{P})\}.
%\]
%
%\begin{thm}[Orthogonal Projection]
%\[
%\mathbb{E}[X \mid \sigma(g)] \text{ is the orthogonal projection of } X \text{ onto } L^2(\sigma(g)).
%\]
%\end{thm}
%
%\begin{proof}
%For $Z \in L^2(\sigma(g))$, we have
%\[
%\langle X - \mathbb{E}[X\mid \sigma(g)], Z \rangle = 0,
%\]
%by the defining property of conditional expectation:
%\[
%\int_A (X - \mathbb{E}[X \mid \sigma(g)]) \, d\mathbb{P} = 0, \quad \forall A \in \sigma(g).
%\]
%\end{proof}
%
%\begin{center}
%\begin{tikzpicture}[scale=3]
%\draw[thick] (0,0) -- (1,0) node[right]{$L^2$};
%\draw[thick] (0,0) -- (0,1) node[above]{$X$};
%\draw[blue,thick] (0,0) -- (0.7,0.7) node[right]{$L^2(\sigma(g))$};
%\draw[red,thick,dashed] (0.7,0.7) -- (0.7,0.3) node[right]{$\mathbb{E}[X|\sigma(g)]$};
%\end{tikzpicture}
%\captionof{figure}{Geometric projection of $X$ onto $L^2(\sigma(g))$}
%\end{center}
%
%\section{Examples with Diagrams}
%
%\subsection*{Example 1: $g(x) = x^2$ on $[0,1]$}
%
%Fibers: $\{x, 1-x\}$.
%\[
%\mathbb{E}[\sin(5\pi x) \mid \sigma(x^2)] = \frac12(\sin(5\pi x) + \sin(5\pi (1-x))).
%\]
%
%\begin{center}
%\begin{tikzpicture}[scale=4]
%\draw[->] (0,0) -- (1,0) node[right]{$x$};
%\draw[->] (0,0) -- (0,1) node[above]{$x^2$};
%\draw[blue, thick] (0,0) parabola (1,1);
%\draw[red, dashed] (0.8,0) -- (0.64,0.64) -- (0.2,0.04) -- cycle;
%\end{tikzpicture}
%\captionof{figure}{Fibers of $g(x)=x^2$; each horizontal slice shows points with same $x^2$}
%\end{center}
%
%\subsection*{Example 2: $g(x)=|x|$ on $[-1,1]$}
%
%Fibers: $\{-y,y\}$.
%\[
%\mathbb{E}[f(x) \mid \sigma(|x|)] = \frac{f(x) + f(-x)}{2}.
%\]
%
%\begin{center}
%\begin{tikzpicture}[scale=3]
%\draw[->] (-1,0) -- (1,0) node[right]{$x$};
%\draw[->] (0,0) -- (0,1) node[above]{$|x|$};
%\draw[blue, thick] (-1,1) -- (0,0) -- (1,1);
%\draw[red, dashed] (-0.8,0.8) -- (0.8,0.8);
%\end{tikzpicture}
%\captionof{figure}{Fibers of $g(x)=|x|$; symmetric points are identified}
%\end{center}
%
%\subsection*{Example 3: Dyadic $\sigma$-Algebra}
%
%Partition $[0,1]$ into $2^n$ intervals.
%For $f(x)=\sin(40\pi x)$:
%\[
%\mathbb{E}[f \mid \mathcal{G}_n](x) = 2^n \int_{I_k} f(t) \, dt, \quad x \in I_k.
%\]
%
%\begin{center}
%\begin{tikzpicture}[scale=4]
%\draw[->] (0,0) -- (1,0) node[right]{$x$};
%\foreach \k in {0,0.25,0.5,0.75,1} {
%    \draw[dashed] (\k,0) -- (\k,0.5);
%}
%\draw[blue, thick] plot[smooth] coordinates {(0,0) (0.125,0.1) (0.25,0.5) (0.375,0.3) (0.5,0.4) (0.625,0.2) (0.75,0.5) (0.875,0.35) (1,0.1)};
%\end{tikzpicture}
%\captionof{figure}{Piecewise averaging over dyadic intervals}
%\end{center}
%
%\section*{Summary}
%
%\begin{itemize}
%    \item $\sigma(g)$ is generated by preimages of measurable sets, equivalently unions of fibers.
%    \item Functions measurable w.r.t. $\sigma(g)$ are exactly $h \circ g$.
%    \item Conditional expectation $\mathbb{E}[X\mid \sigma(g)]$ is the Radon--Nikodym derivative of $\nu(A) = \int_A X \, d\mathbb{P}$ on $\sigma(g)$.
%    \item In $L^2$, $\mathbb{E}[X \mid \sigma(g)]$ is the orthogonal projection onto $L^2(\sigma(g))$.
%    \item Examples illustrate averaging over fibers, symmetric functions, and dyadic projections.
%\end{itemize}
%
%
%
%\section{Conditional Expectation in $L^2([0,1]^2)$}
%
%Let $(X,Y)$ be random variables on $[0,1]^2$ with Lebesgue measure.
%Let $f \in L^2([0,1]^2)$, and let $g(X,Y)$ be a measurable function.
%The conditional expectation $\mathbb{E}[f \mid \sigma(g)]$ is the orthogonal projection of $f$ onto the subspace
%\[
%L^2(\sigma(g)) = \{ h \circ g : h \in L^2(S, \mathcal{B}, g_\#\mathbb{P})\}.
%\]
%
%\subsection{Example 1: $g(x,y) = x+y$}
%
%The fibers are lines:
%\[
%\{(x,y) \in [0,1]^2 : x+y = c\}.
%\]
%
%Conditional expectation averages $f$ along lines of constant sum.
%
%\begin{center}
%\begin{tikzpicture}[scale=5]
%% axes
%\draw[->] (0,0) -- (1.1,0) node[right]{$x$};
%\draw[->] (0,0) -- (0,1.1) node[above]{$y$};
%
%% fibers x+y = c
%\foreach \c in {0.2,0.4,0.6,0.8,1.0,1.2} {
%    \draw[red, thick, dashed] (0,\c) -- (\c,0);
%}
%
%\node at (0.9,0.1) {fibers $x+y=c$};
%\end{tikzpicture}
%\captionof{figure}{Fibers of $g(x,y)=x+y$ in $[0,1]^2$; conditional expectation averages along these lines.}
%\end{center}
%
%\subsection{Example 2: $g(x,y) = \lfloor 2x \rfloor + \lfloor 2y \rfloor$}
%
%Partition $[0,1]^2$ into four $2 \times 2$ blocks:
%\[
%[0,0.5)\times[0,0.5),\quad [0.5,1]\times[0,0.5),\quad [0,0.5)\times[0.5,1],\quad [0.5,1]\times[0.5,1].
%\]
%
%Conditional expectation is constant on each block.
%
%\begin{center}
%\begin{tikzpicture}[scale=4]
%\draw[thick] (0,0) rectangle (1,1);
%\draw[thick] (0.5,0) -- (0.5,1);
%\draw[thick] (0,0.5) -- (1,0.5);
%\foreach \x in {0.25,0.75} {
%    \foreach \y in {0.25,0.75} {
%        \fill[blue!20] (\x-0.25,\y-0.25) rectangle (\x+0.25,\y+0.25);
%    }
%}
%\node at (0.25,0.25) {Block 1};
%\node at (0.75,0.25) {Block 2};
%\node at (0.25,0.75) {Block 3};
%\node at (0.75,0.75) {Block 4};
%\end{tikzpicture}
%\captionof{figure}{Conditional expectation constant on each dyadic block for $g(x,y)=\lfloor 2x \rfloor + \lfloor 2y \rfloor$.}
%\end{center}
%
%\section{Geometric Projection in $L^2([0,1]^2)$}
%
%Let $f \in L^2([0,1]^2)$ and $H = L^2(\sigma(g))$. Then $\mathbb{E}[f \mid \sigma(g)]$ is the unique function $h \in H$ such that
%\[
%\langle f - h, Z \rangle = 0, \quad \forall Z \in H,
%\]
%i.e.,
%\[
%\iint_{[0,1]^2} (f(x,y) - \mathbb{E}[f \mid \sigma(g)](x,y)) Z(x,y) \, dx dy = 0.
%\]
%
%\begin{center}
%\begin{tikzpicture}[scale=5]
%% axes
%\draw[->] (0,0) -- (1.1,0) node[right]{$x$};
%\draw[->] (0,0) -- (0,1.1) node[above]{$y$};
%
%% subspace L^2(sigma(g))
%\draw[blue, thick] (0,0) -- (0.9,0.9) node[right]{$L^2(\sigma(g))$};
%
%% vector X
%\draw[->, red, thick] (0,0) -- (0.6,0.8) node[above right]{$f$};
%
%% projection
%\draw[dashed, red] (0.6,0.8) -- (0.7,0.7) node[right]{$\mathbb{E}[f \mid \sigma(g)]$};
%\end{tikzpicture}
%\captionof{figure}{Projection of $f$ onto the subspace of $\sigma(g)$-measurable functions in $L^2([0,1]^2)$.}
%\end{center}
%
%\section{Summary}
%
%\begin{itemize}
%    \item In multidimensional $L^2$, fibers are level sets of $g$ in $[0,1]^d$.
%    \item Conditional expectation averages $f$ along fibers, giving the orthogonal projection onto $L^2(\sigma(g))$.
%    \item Examples include sums, floor partitions, and dyadic blocks.
%    \item Graphical representation helps visualize how conditional expectation reduces information along fibers.
%\end{itemize}
%
%
%
%\section{Circular Fibers Example: $g(x,y)=x^2+y^2$}
%
%Consider $f \in L^2([0,1]^2)$ and $g(x,y) = x^2 + y^2$.
%The fibers are circles (or portions of circles) with constant radius:
%\[
%\{(x,y) \in [0,1]^2 : x^2 + y^2 = r^2\}.
%\]
%
%Conditional expectation $\mathbb{E}[f \mid \sigma(g)]$ averages $f$ along each circular fiber.
%
%\subsection{3D Perspective of Fibers}
%
%\tdplotsetmaincoords{60}{120} % set view angle
%\begin{center}
%\begin{tikzpicture}[tdplot_main_coords, scale=4]
%
%% axes
%\draw[->] (0,0,0) -- (1.2,0,0) node[below right]{$x$};
%\draw[->] (0,0,0) -- (0,1.2,0) node[above left]{$y$};
%\draw[->] (0,0,0) -- (0,0,1.2) node[above]{$g(x,y)$};
%
%% grid plane
%\draw[gray!50, dashed] (0,0,0) grid (1,1,0);
%
%% circular fibers at various heights
%\foreach \r in {0.2,0.4,0.6,0.8,1.0} {
%    \draw[red, thick, dashed]
%        ({\r*cos(45)},{\r*sin(45)},{\r^2});
%    \draw[red, thick, dashed]
%        ({\r*cos(135)},{\r*sin(135)},{\r^2});
%    \draw[red, thick, dashed]
%        ({\r*cos(225)},{\r*sin(225)},{\r^2});
%    \draw[red, thick, dashed]
%        ({\r*cos(315)},{\r*sin(315)},{\r^2});
%}
%
%\node at (0.5,0.5,1) {Fibers $x^2+y^2=r^2$};
%
%\end{tikzpicture}
%\captionof{figure}{3D perspective of circular fibers for $g(x,y)=x^2+y^2$. Conditional expectation averages $f$ along circles.}
%\end{center}
%
%\subsection{Conditional Expectation Averaging}
%
%For any $(x,y)$, the conditional expectation is
%\[
%\mathbb{E}[f \mid \sigma(g)](x,y) = \frac{1}{\mu(\text{fiber})} \int_{\text{fiber of }(x,y)} f(u,v) \, d\mu(u,v),
%\]
%where $\mu$ is the 1D measure along the circular fiber.
%
%\subsection{Geometric Projection in $L^2([0,1]^2)$}
%
%\begin{center}
%\begin{tikzpicture}[tdplot_main_coords, scale=4]
%
%% axes
%\draw[->] (0,0,0) -- (1.2,0,0) node[right]{$x$};
%\draw[->] (0,0,0) -- (0,1.2,0) node[left]{$y$};
%\draw[->] (0,0,0) -- (0,0,1.2) node[above]{$f(x,y)$};
%
%% vector f
%\draw[->, red, thick] (0,0,0) -- (0.7,0.7,0.9) node[right]{$f$};
%
%% projection onto fiber-constant subspace
%\draw[dashed, blue, thick] (0,0,0) -- (0.7,0.7,0.5) node[right]{$\mathbb{E}[f|\sigma(g)]$};
%
%\draw[blue, thick] (0.7,0.7,0) -- (0.7,0.7,0.5);
%
%\end{tikzpicture}
%\captionof{figure}{Projection of $f$ onto the subspace of circular-fiber-constant functions in $L^2([0,1]^2)$.}
%\end{center}
%
%\section{Summary}
%
%\begin{itemize}
%    \item Circular fibers $x^2+y^2=r^2$ are level sets of $g$.
%    \item Conditional expectation averages $f$ along these fibers.
%    \item In $L^2$, this corresponds to orthogonal projection onto the subspace of functions constant along fibers.
%    \item The 3D diagrams illustrate the “fiber-constant” subspace and the geometric projection of $f$.
%\end{itemize}



{\bf Declarations}
\begin{itemize}
\item Conflicts of Interest:
The authors declare no conflicts of interest.
\item Author Contributions:
The authors have the same contribution.
\item Data Availability Statement:
Data is not applicable. The results are obtained by manual computations.
\item Funding:
No funding.
\end{itemize}
\begin{thebibliography}{99}

%\bibitem{as1} J. Agler and M. Stankus, m-Isometric transformations of Hilbert space, I. Integr equ oper theory {\bf 21} (1995) 383-–429.
%\bibitem{as2} J. Agler and M. Stankus, m-Isometric transformations of Hilbert space, II. Integr equ oper theory {\bf 21} (1995) 1-–48.
%\bibitem{as3} J. Agler and M. Stankus, m-Isometric Transformations of Hilbert Space, III, Integr equ oper theory {\bf 24} (1996) 379-–421.
%\bibitem{ber} T. Bermudez, A. Martinon and J.A. Noda, Products of m-isometries. Linear Algebra and its Applications {\bf 438} (2013) 80-–86.
%\bibitem{bm} T. Bermudez, I. Marrero and A. Martinon, On the orbit of an m-isometry. Integral Equations and Operator Theory {\bf 64} (2009) 487-–494.
%\bibitem{bj} F. Botelho, and J. Jamison,  Isometric properties of elementary operators. Linear Algebra and its Applications {\bf 432} (2010) 357–-365.
%\bibitem{bu} M. Buchała, m-Isometric Composition Operators on Discrete Spaces. Complex Anal. Oper. Theory, {\bf 158} (2024).
%\bibitem{bbl} C. Burnap, I. Bong Jung and A. Lambert, Separating partial normality classes with composition operators, J. Operator theory, {\bf 53} (2005) 381--397.
%\bibitem{cj} J. T. Campbell and J. E. Jamison, On some classes of weighted composition operators, Glasgow Mathematical Journal, {\bf32} (1990) 87--94.
%\bibitem{cot} M. Cho, S. Ota and K. Tanahashi, Invertible weighted shift operators which are m-isometries.
%Proceedings of the American Mathematical Society {\bf141} (2013) 4241–-4247.
%\bibitem{Douglas1966}
%R. G. Douglas, On majorization, factorization, and range inclusion of operators on Hilbert space. Proceedings of the American Mathematical Society, \textbf{17}(2), (1966), 413–415.
\bibitem{dhd} P.G. Dodds, C.B. Huijsmans and B. De Pagter,
characterizations of conditional expectation-type operators,
Pacific J. Math. {\bf 141}(1) (1990), 55-77.
\bibitem{die} J. Diestel, H. Jarchow and A. Tonge, Absolutely summing operators, Cambridge
Studies in Advanced Mathematics, vol. 43, Cambridge University Press, Cambridge, (1995).
MR1342297.
\bibitem{emj} H. Emamalipour and M. R. Jabbarzadeh, Lambert Conditional Operators on $L^2(\Sigma)$. Complex Anal. Oper. Theory 14, 18 (2020).
\bibitem{e2} Y. Estaremi, Multiplication conditional expectation type operators on Orlicz spaces, Journal of Mathematical Analysis and Applications {\bf414} 88-98.
\bibitem{e} Y. Estaremi, On a Class of Operators With Normal Aluthge Transformations, Filomat {\bf 29} (2015) 969-–975.
\bibitem{e1} Y. Estaremi, On the algebra of WCE operators, Rocky Mountain J. Math. {\bf 48} (2018) 501-- 517.
\bibitem{ej} Y. Estaremi and M.R. Jabbarzadeh, Weighted lambert type operators on
$L^{p}$-spaces, Oper. Matrices {\bf 1} (2013), 101-116.
\bibitem{ej2} Y. Estaremi and M.R. Jabbarzadeh, Compact Lambert type operators between two Lp spaces, J. Math. Anal. Appl. {\bf 420} (2014) 118–-123.
%$L^{p}$-spaces, to appear in Mediterranean Journal of Mathematics.
%\bibitem{th} T. Huruya, A note on $p$-hyponormal operators, Proc.
%Amer. Math. Soc. {\bf 12} (1997), 3617-3624.
%\bibitem{jl} S. Jung, and J.E. Lee, On (A;m)-expansive operators. Studia Mathematica {\bf213} (2012) 3-–23.
%\bibitem{lam} A. Lambert, Hyponormal composition operators. Bull. Lond. Math. Soc. {\bf 18} (1986) 395-–400.
\bibitem{her}
J. Herron, Weighted conditional expectation operators, Oper.
Matrices {\bf 1} (2011), 107-118.
%\bibitem{mp} S. Mecheri and T. Prasad, On n-quasi-m-isometric operators. Asian European Journal of
%Mathematics {\bf9} (2016).
%\bibitem{mp1} S. Mecheri and S. M. Patel, On quasi-2-isometric operators. Linear Multi-Linear Algebra, (2017) 1019–-1025. https://doi.org/10.1080/03081087.2017.1335283.
%%\bibitem{mo}
%Shu-Teh Chen, Moy,  Characterizations of conditional expectation
%as a transformation on function spaces, Pacific J. Math. {\bf 4}
%(1954), 47-63.
%
%
%\bibitem{rao1}
%M. M. Rao, Conditional measure and operators, J. Multivariate
%Anal. 5 (1975) 330-413.

%\bibitem{rao2}
%M. M. Rao, two characterizations of conditional probability, Proc.
%Amer. Math. Soc. 19 (1976) 75-80.
\bibitem{pi}
A.~Pietsch,
\emph{Nuclear Locally Convex Spaces},
Ergebnisse der Mathematik und ihrer Grenzgebiete, Band~66,
Springer-Verlag, New York--Heidelberg, 1972.

\bibitem{rao}
M. M. Rao, Conditional measure and applications, Marcel Dekker,
New York, 1993.
\bibitem{ri1}
O. I. Reinov, Approximation Properties and Some Classes of Operators. Journal of Mathematical Sciences {\bf 107} (2001) 3911-–3951.
\bibitem{ri}
O. I. Reinov, On linear operators with s-nuclear adjoints, $0<s\leq1$, J. Math. Anal. Appl. {\bf 415} (2014) 816-–824.
\bibitem{z}
A. C. Zaanen, Integration, 2nd ed., North-Holland, Amsterdam,
1967.

\end{thebibliography}
\end{document}
